La caractéristique intensité-tension d’un générateur est la courbe
$U_{PN}=f(I)$ représentée sur le schéma ci-dessous.

\begin{figure}[H]
    \centering
    \begin{tikzpicture}
        \begin{axis}[
                cycle list name=my color list,
                width=9cm,height=6cm,
                xlabel=$I~(\unit{\A})$,ylabel=$U~(\unit{\volt})$,
                xmin=0,xmax=0.95,
                ymin=0,ymax=5.5,
                xtick distance=0.1,
                ytick distance=1,
            ]
            \addplot {-1.5/0.7*\x+4.5};
        \end{axis}
    \end{tikzpicture}
\end{figure}

\begin{enumerate}
    \item Déterminer la fém $E$ et la résistance interne $r$ de ce
          générateur.
    \item Calculer la puissance électrique qu'il fournit au circuit quand le
          courant qu'il débite a une intensité $I=\qty{0,25}{\ampere}$.
          Quelle est l'énergie totale convertie si le générateur fonctionne
          pendant 20 minutes~?
\end{enumerate}


